\documentclass[10pt,a4paper]{amsart}
\usepackage[margin=19mm]{geometry}
\usepackage{amsmath,amssymb,mathtools}
\usepackage[scr=rsfs,cal=euler]{mathalfa}
\usepackage{xfrac}
\usepackage[unicode,hidelinks,bookmarks=false]{hyperref}
\newcommand{\ie}{i.e.\ }
\newcommand{\cf}{cf.\ }
\newcommand{\resp}{respectively\ }
\newcommand{\Subsection}{\subsection}
\providecommand{\nobreakdash}{\nobreak\hbox{-}}
\setlength{\emergencystretch}{2em}
\multlinegap0pt
\usepackage{mathtools}
\usepackage{amssymb}
\usepackage{cancel}
\usepackage{bm}
\newtheorem{proposition}{Proposition}
\newtheorem{lemma}{Lemma}
\newcommand{\Ib}{\cI_\bullet}
\newcommand{\Sb}{\bSig_\bullet}
\DeclareMathOperator{\Supp}{Supp}
\newcommand{\bSig}{\mathbf\Sigma}
\newcommand{\bi}{\mathbf{i}}
\newcommand{\bt}{\mathbf{t}}
\newcommand{\cC}{\mathcal{C}}
\newcommand{\cI}{\mathcal{I}}
\newcommand{\cT}{\mathcal{T}}
\newcommand{\sbb}{\sigma_\bullet}
\DeclareMathOperator{\sgn}{sgn}
\title{Construction of Curves with a Controlled First Slope using p-Symmetric Numbers}
\author{Robert Moore, Hui June Zhu}
\date{Source: arXiv:2411.01832v3 (CC BY 4.0)}
\begin{document}
\maketitle

\noindent\emph{Source and licence.} Construction of Curves with a Controlled First Slope using p-Symmetric Numbers. Version arXiv:2411.01832v3 (CC BY 4.0), distributed by arXiv under the Creative Commons Attribution 4.0 International licence, http://creativecommons.org/licenses/by/4.0/, as recorded in the arXiv licence metadata for this exact version and shown on the abstract page https://arxiv.org/abs/2411.01832v3. Full licence notice: \texttt{licenses/arxiv-2411.01832-CC-BY-4.0.txt}.\par\smallskip

\noindent\textit{Editorial scope.} Counted unit: the article's proposition P:bound (source0.tex lines 1362--1390). The article's complete original proof of it is delivered with it (source0.tex lines 1393--1458, one \texttt{proof} environment of 66 lines). No mathematics is written, repaired or re-derived: the statement and the proof are the article's own lines, in the article's own notation, and no retained line is substituted or re-flowed.  Retained uncounted as the source-local context the proof needs: lines 1026--1035; lines 1075--1086; lines 1174--1176; lines 1189--1212; lines 1259--1273; lines 1355--1358.  Nothing was omitted from any retained span, so the package carries no omission record.  No retained line contains a citation, so the wrapper carries no bibliography and no bibliography entry.  No retained line contains a figure environment, an image-inclusion command or drawing code, so no image is delivered and the package declares no asset dependency.  Editorial formatting only, in addition: an AMS \texttt{amsart} preamble that restates the article's own declarations and macros used by the retained lines, namely the environments \texttt{proposition}, \texttt{lemma} on separate counters, together with the standard packages those lines need.  The retained environments print in this document as Proposition 1, Lemma 1 in order of appearance, whereas the article prints the same items with the numbering of its own full document.  The complete native source of the article as distributed in the arXiv e-print for this exact version is bundled unchanged as \texttt{sources/167913/source0.tex}.
\begin{proposition}
\label{P:uniq-perm}
Let $\cC=\cC(\bi,a)$ where $\bi$ contains a unique element $\nu$ with maximal $p$-adic weight. 
\begin{enumerate}
\item If $\nu$ is $p$-symmetric with minimal factorization $\nu * w=(p^k-1)\ell$
and $k|a$, then $\Sb\neq \varnothing$. 
\item Suppose $\nu<p^a$. If $\Sb\neq \varnothing$, 
then $\nu$ is $p$-symmetric. 
\end{enumerate}
\end{proposition}

\begin{proposition}
\label{P:bounded_min2}
Let $\cC$ be as in Proposition \ref{P:uniq-perm}.
\begin{enumerate}
\item
If $\nu$ is $p$-symmetric with minimal factorization 
$\nu * w = (p^k-1)\ell$ where $k|a$ and has shift factor $e$,  
then $k-e\le t(\cC)\le \nu$. 
\item 
If $\nu=p^k-1$, $k|a$, then $t(\cC)=k$.
\end{enumerate}
\end{proposition}

\begin{equation}\label{E:coin-set2}
\fbox{$\cC=\cC(\Supp(f),a)=\{ip^j\mid (i,j)\in I\}=\{ip^j\mid i\in\Supp(f), j=0,\ldots,a-1\}.$}
\end{equation}

\begin{lemma}
\label{L:G_N1}
Write
$E(x)=\sum_{n=0}^{\infty} \beta_n x^n$ and 
$F_0(x)=\sum_{N=0}^{\infty} G_N x^N$.
\begin{enumerate}
\item Then 
$
G_N = \sum_{\bt \in \cT(N)}\kappa(\bt) \gamma^{|\bt|},
$
where 
$
\kappa(\bt)= \prod_{(i,j)\in I}\beta_{t_{ij}}\hat a_i^{p^j t_{ij}}.
$
\item Suppose $\nu$ is a maximal $p$-adic weight element in $\Supp(f)$.
Then 
$$v_p(G_N)
\ge \frac{M_\cC(N)}{p-1}
\ge \frac{s_p(N)}{(p-1)s_p(\nu)}.$$ 
If $\nu$ is the unique maximal $p$-adic weight element and the second inequality is an equality, then 
$v_p(G_N)=\frac{s_p(N)}{(p-1)s_p(\nu)}$.
\end{enumerate}

\end{lemma}

\begin{lemma}
\label{L:G_N}
Let $G_N$ be as in Lemma \ref{L:G_N1}.
Suppose $\nu$ is a maximal $p$-adic weight element in $\Supp(f)$.
For any $m\ge 1$, if $(\cI,\sigma)\in \bSig_m$, then  
    $
        v_p\left(\prod_{\ell\in \cI} G_{q\ell - \sigma \ell}
            \right) \ge \frac{am}{s_p(\nu)}.
    $
If $\nu$ is the unique maximal $p$-adic weight element in $\Supp(f)$ and
$\nu<q$, then 
$(\cI,\sigma)$ is a minimizer if and only if  
$v_p\left(\prod_{\ell\in \cI} G_{q\ell - \sigma \ell}
            \right) = \frac{am}{s_p(\nu)}$.
\end{lemma}

\begin{equation}
\label{E:C_m} 
C_m=\sum_{(\cI,\sigma)\in \bSig_m}(-1)^m\sgn(\sigma)\prod_{\ell\in \cI}G_{q\ell-\sigma\ell}.   
\end{equation}

\begin{proposition}\label{P:bound}
Let $\cC=\cC(\Supp(f),a)$ be the coin set as in \eqref{E:coin-set2}. 
\begin{enumerate}
\item 
Let $\nu$ be a maximal $p$-adic weight element in $\Supp(f)$. Then     
$
        v_p(C_m) \ge \frac{am}{s_p(\nu)}
    $ for all $m\ge 1$.   
\item 
Suppose $\nu$ is the unique maximal $p$-adic weight element in $\Supp(f)$ and $\nu<q$.
\begin{enumerate}
\item If $\nu$ is $p$-symmetric with minimal factorization
    $\nu * w = (p^k-1)\ell$ and $k|a$,
then $\Sb\neq \varnothing$.
Let $t \coloneqq t(\cC)$
be the minimizer height of $\cC$,
and let $e$ denote the shift factor of $\nu$.
Then we have $k-e\le t\le \nu$, $v_p(C_t)=\frac{at}{s_p(\nu)}$, and 
$v_p(C_m)>\frac{am}{s_p(\nu)}$ for all $m>t$.
\item 
Conversely, if $v_p(C_m)=\frac{am}{s_p(\nu)}$ for some $m\ge 1$, then 
$\nu$ is $p$-symmetric.    
\end{enumerate}
\item Suppose $p^k-1$ with $k|a$ is the unique maximal $p$-adic weight element in $\Supp(f)$.
Then $v_p(C_m) \ge \frac{am}{k(p-1)}$ for all $m\ge 1$,
$v_p(C_k)=\frac{a}{p-1}$, 
and $v_p(C_m)>\frac{am}{k(p-1)}$ for all $m>k$.
\end{enumerate}
\end{proposition}

\begin{proof}
(1) By Lemma \ref{L:G_N},
$ v_p(\prod_{\ell\in \cI}
    G_{q\ell - \sigma \ell})\ge \frac{am}{s_p(\nu)}
$ for every $(\cI,\sigma)\in \bSig_m$.
So their sum still has 
$v_p(C_m) \ge \min_{(\cI,\sigma)\in\bSig_m} v_p(\prod_{\ell\in \cI}
    G_{q\ell - \sigma \ell})\ge \frac{am}{s_p(\nu)}.$
\\
(2a)
$\Sb \ne \varnothing$ immediately follows from
Proposition \ref{P:uniq-perm}(1),
and the inequality $k-e \le t \le \nu$ immediately follows from
Proposition \ref{P:bounded_min2}(1).

By \eqref{E:C_m}, we have 
$
C_t=C_{t,-} + C_{t,+}
$
where 
\begin{equation*}
    C_{t,+}= 
    \sum_{\stackrel{(\cI,\sigma)\in \bSig_t}{(\cI,\sigma)\ne (\Ib,\sbb)}}
    (\pm)\prod_{\ell\in \cI} G_{q \ell - \sigma\ell},\qquad
C_{t,-}=\pm \prod_{\ell\in \Ib}G_{q\ell-\sbb\ell}.
\end{equation*}
To prove that $v_p(C_t)=\frac{at}{s_p(\nu)}$, it suffices to show
the inequality
$v_p(C_{t,+})>\frac{at}{s_p(\nu)}$
and the equality
$v_p(C_{t,-})=\frac{at}{s_p(\nu)}$.   

Let  $(\cI,\sigma)\in \Sigma_t$ and $(\cI,\sigma)\neq (\Ib,\sbb)$.
Since $(\Ib,\sbb)$ is the unique maximal minimizer,
$(\cI,\sigma)$ is not a minimizer.
By Lemma \ref{L:G_N}, we have
 $
        v_p\left(\prod_{\ell \in \cI} G_{q\ell - \sigma\ell}\right)
            > \frac{at}{s_p(\nu)}
$,
hence
 $v_p(C_{t,+})>\frac{at}{s_p(\nu)}$.
On the other hand, by Lemma \ref{L:G_N}, since $(\Ib,\sbb)$ is a minimizer,
we have
$v_p(C_{t,-})=
v_p(\prod_{\ell\in \Ib}G_{q\ell-\sbb\ell}) = \frac{at}{s_p(\nu)}$.

Suppose $m>t$.
Then $\bSig_m$ contains no minimizers by the maximality of $(\Ib,\sbb)$, 
hence \eqref{E:C_m} and Lemma \ref{L:G_N} show that 
$v_p(C_m)>\frac{am}{s_p(\nu)}$.

\noindent 
(2b)
Suppose $v_p(C_m)=\frac{a m}{s_p(\nu)}$ for some $m\ge 1$.
From (\ref{E:C_m}) and Lemma \ref{L:G_N},
there must exist a pair $(\cI,\sigma)\in \bSig_m$ 
that is a minimizer.
By  Proposition \ref{P:uniq-perm}(2), since $\nu<q$, 
we conclude that $\nu$ is $p$-symmetric.    
\\
(3) 
Notice that  $\nu=p^k-1$ is $p$-symmetric.
The minimizer height $t(\cC)=k$ by Proposition \ref{P:bounded_min2}.
The statement follows from the above argument. 
\end{proof}

\end{document}
